% =============================================================================
% SECTION 7: THREATS TO VALIDITY
% =============================================================================
\section{Threats to Validity}
\label{sec:threats}

We discuss threats to validity following the Wohlin et al.\ taxonomy of construct,
internal, external, and conclusion validity.

\subsection{Construct Validity}

\paragraph{Classification Subjectivity.}
The primary PAR classification (Field F2 in the data extraction template) requires a judgment
about which paradigm is most central to a given system. For hybrid systems, this judgment is
inherently subjective. We mitigated this threat through inter-rater reliability testing
($\kappa = 0.78$) and by recording secondary PARs explicitly. However, there remains a risk
that some systems were classified under a paradigm that captures only one aspect of their
approach. Sensitivity analyses in the replication package show that reclassifying the 15 most
ambiguous cases does not materially change the findings.

\paragraph{Task Taxonomy Completeness.}
Our seven-task taxonomy covers over 95\% of task targets in the corpus, but tasks in the
remaining 5\% (e.g., code review, documentation generation, architectural recovery) are
excluded from the task--paradigm mapping. These tasks share properties with the included tasks
and would likely follow similar paradigm distributions, but we cannot confirm this without
systematic analysis.

\subsection{Internal Validity}

\paragraph{Selection Bias.}
Despite our systematic search protocol, some relevant papers may have been missed. Authors
who do not use the terms ``repository-level'' or ``codebase'' in their title or abstract may
have been excluded from the initial candidate pool. We partially addressed this through
snowballing, but forward snowballing is limited to papers that cite already-included studies.
Emerging paradigms that have not yet attracted many citations may be under-represented.

\paragraph{Publication Bias.}
As with all literature reviews, positive results are more likely to be published and indexed
than negative results. Studies that found that a particular paradigm \emph{does not} improve
over baselines are less likely to appear in our corpus. This means that our performance
characterization (RQ3) overestimates average paradigm performance relative to what would be
observed in a true random sample of deployment attempts.

\paragraph{Temporal Bias.}
Our search was frozen in January 2026. The field is evolving rapidly, and paradigms that
appear underrepresented in our corpus (e.g., \parcmd{6}, \parcmd{5}) may have received
substantially more attention by the time this paper is published. We note this limitation but
cannot resolve it without continuous updating, which is beyond the scope of a single SLR.

\subsection{External Validity}

\paragraph{Generalizability to Proprietary Systems.}
Our corpus is exclusively composed of published research papers; proprietary systems deployed
by large software companies (GitHub Copilot Enterprise, JetBrains AI Assistant, Microsoft
IntelliCode) use repository representation techniques that are not disclosed in the literature.
Our findings may not generalize to these systems, which may employ substantially more
sophisticated approaches.

\paragraph{Programming Language Distribution.}
As noted in Section~\ref{subsec:rq5}, approximately 68\% of studies evaluate on Python
repositories. Findings about paradigm effectiveness may not generalize to statically typed
languages (Java, C\#, TypeScript) where type information plays a larger role, or to systems
languages (C, Rust) where memory safety and low-level resource management impose additional
constraints.

\subsection{Conclusion Validity}

\paragraph{Metric Heterogeneity.}
The diverse set of evaluation metrics used across studies (Exact Match, Edit Similarity,
CodeBLEU, pass@k, MRR, Top-k accuracy, F1) makes cross-study comparison difficult.
Table~\ref{tab:completion-perf} and Table~\ref{tab:swebench-perf} are restricted to the
two benchmarks with the most standardized metrics. For other tasks and benchmarks, we report
findings at the qualitative level rather than the quantitative level to avoid misleading
comparisons.

\paragraph{Non-Reproducibility.}
A significant fraction of studied systems (approximately 22\%) do not release source code or
model weights, making it impossible to independently verify reported results or to run
controlled comparisons. This limits the precision of our performance characterization and is
itself an important finding: the software engineering research community should strengthen
norms around artifact availability.
